\section{Algorithm Design under a Fixed Feature Set}
\label{sec:algorithms}

Sections~\ref{sec:marginal} and~\ref{sec:linvsnonlin} varied the information set
and the model class. This section holds both fixed. The feature construction is
frozen at the pipeline of Section~\ref{sec:descriptive}, as are the
$n = 218{,}934$ out-of-sample panel, the evaluation protocol (QLIKE on raw
variance after Duan smearing --- a loss robust to noise in the variance proxy in
the sense of \citealp{Patton2011,HansenLundeNason2011}; Diebold--Mariano with HAC
standard errors and the Harvey--Leybourne--Newbold correction,
\citealp{DieboldMariano1995,HarveyLeybourneNewbold1997}), and the incumbent:
per-bar ordinary least squares on HAR plus calendar, QLIKE $0.13415$. What
remains free is the \emph{algorithm}.

\subsection{The design space}
\label{sec:alg_principle}

With the design matrix frozen, four levers remain, and this section presents one
algorithm per lever. \textbf{(i) Exactness and cadence:} a walk-forward backtest
slides its window one bar at a time, so the estimator can be refit from scratch
(exact, expensive), refit on a cadence (cheap, stale), or \emph{updated} exactly as
the window slides. \textbf{(ii) Causal hyperparameter choice:} a hyperparameter
selected on the whole sample is a fitted value, not a forecast; choosing it only
from data strictly preceding each forecast turns ``tuning does not help'' into a
claim about the data rather than about the analyst's discipline.
\textbf{(iii) Composition:} in a base-plus-correction stack the base does not
depend on the correction's configuration, and that redundancy can be removed
exactly. \textbf{(iv) Geometry:} all of the above are parametric least squares; the
alternative is to ask which past episodes \emph{resemble} the present one, in a
metric learned from the data. Each subsection gives an idea, the device that makes
it feasible under per-bar refitting, and a measured verdict.

\subsection{Exact every-bar linear updates}
\label{sec:alg_exact}

\subsubsection{Rank-one rolling least squares}

Let the training window at bar $t$ be the $W$ rows $[t-W, t)$ with design $X_t$
and target $y_t$. Ridge regression with an unpenalized intercept
\citep{HoerlKennard1970} solves the centred normal equations
\begin{equation}
\bigl( X_t^{c\top} X_t^{c} + \alpha I_p \bigr)\,\beta_t = X_t^{c\top} y_t^{c},
\qquad X_t^{c} = X_t - \bar{x}_t \mathbf{1}^{\top},
\label{eq:ridge_normal}
\end{equation}
which costs $O(Wp^2 + p^3)$ per refit. But the window changes by one row in and
one row out, so its sufficient statistics change by rank one. Maintaining
$S^{xx}_t = X_t^\top X_t$, $S^{xy}_t = X_t^\top y_t$, $s^{x}_t = X_t^\top
\mathbf{1}$ and $s^{y}_t = y_t^\top \mathbf{1}$, the slide that admits
$(x_{\mathrm{in}}, y_{\mathrm{in}})$ and evicts $(x_{\mathrm{out}},
y_{\mathrm{out}})$ is
\begin{equation}
\begin{aligned}
S^{xx}_{t+1} &= S^{xx}_t + x_{\mathrm{in}} x_{\mathrm{in}}^\top - x_{\mathrm{out}} x_{\mathrm{out}}^\top,
&\quad
S^{xy}_{t+1} &= S^{xy}_t + x_{\mathrm{in}} y_{\mathrm{in}} - x_{\mathrm{out}} y_{\mathrm{out}}, \\
s^{x}_{t+1} &= s^{x}_t + x_{\mathrm{in}} - x_{\mathrm{out}},
&\quad
s^{y}_{t+1} &= s^{y}_t + y_{\mathrm{in}} - y_{\mathrm{out}},
\end{aligned}
\label{eq:rank1_update}
\end{equation}
two rank-one updates at $O(p^2)$ in place of an $O(Wp^2)$ rebuild. Centring is
itself incremental --- the centred Gram and cross-product are $S^{xx} - s^{x}
s^{x\top}/W$ and $S^{xy} - s^{x} s^{y}/W$ --- so the solve reproduces
\texttt{sklearn}'s \texttt{fit\_intercept=True} convention exactly. The $O(p^3)$
solve is taken on a cadence; the updates are taken every bar regardless, so the
state is never stale even when the coefficients are refreshed coarsely. Two
properties matter downstream. At $\alpha = 0$, equation~\eqref{eq:ridge_normal} is
exact ordinary least squares, so this solver \emph{is} how the paper's incumbent is
computed. And the speed-up is measured: against a per-window \texttt{sklearn}
\texttt{Ridge} refit the rolling solver agrees to $1.8 \times 10^{-11}$ and runs
$108\times$ faster, $4.589$ ms/bar reduced to $0.042$ ms/bar --- which is what
makes the per-bar-refit convention used throughout this paper affordable.

For the penalized variant the slide can be pushed onto the inverse itself. With
$K_t = (X_t^\top X_t + \alpha D)^{-1}$ for a diagonal penalty $D$ that is zero on
the intercept, Sherman--Morrison gives a closed-form two-step update --- admit the
entering row $u$, then evict the leaving row $v$:
\begin{equation}
K \;\leftarrow\; K - \frac{(Ku)(Ku)^\top}{1 + u^\top K u},
\qquad\text{then}\qquad
K \;\leftarrow\; K + \frac{(Kv)(Kv)^\top}{1 - v^\top K v},
\label{eq:sherman_morrison}
\end{equation}
with $c \leftarrow c + u y_{\mathrm{in}} - v y_{\mathrm{out}}$ and $\beta = Kc$
available at $O(p^2)$ with no factorization at all.

\subsubsection{The \texorpdfstring{$L_1$}{L1} analogue: recursive elastic net by online homotopy}
\label{sec:alg_homotopy}

Equation~\eqref{eq:sherman_morrison} exists because the ridge solution is a smooth
function of the data. The lasso \citep{Tibshirani1996} and elastic net
\citep{ZouHastie2005} are not: their solutions are piecewise linear with kinks
where the active set changes. An exact every-bar update nevertheless survives, in
piecewise form. The starting point is that an elastic net is a lasso on a ridged
Gram matrix:
\begin{equation}
\min_{b} \; \tfrac12 \|y - Xb\|_2^2 + \mu \|b\|_1 + \tfrac12 \lambda_2 \|b\|_2^2
\;\;\equiv\;\;
\min_{b} \; \tfrac12 b^\top G b - c^\top b + \mu \|b\|_1,
\quad
\begin{aligned}
G &= X^\top X + \lambda_2 I, \\ c &= X^\top y.
\end{aligned}
\label{eq:enet_gram}
\end{equation}
This is more than notational tidiness: the $L_2$ term keeps $G$ strictly positive
definite, so the active-set subproblems stay well posed even under near
collinearity --- the failure mode that sinks explicit-inverse formulations of the
same recursion. The implementation matches \texttt{sklearn}'s
\texttt{ElasticNet(alpha, l1\_ratio, fit\_intercept=False)} over $N$ rows through
$\mu = N \alpha \rho$ and $\lambda_2 = N \alpha (1-\rho)$, $\rho =
\texttt{l1\_ratio}$. The solution is characterized by KKT conditions on the
residual correlation $r = c - Gb$: $r_j = \mu_j s_j$ on active coordinates with
sign $s_j$, and $|r_j| \le \mu_j$ elsewhere. Both the classical $\mu$-path
\citep{Efron2004} and the \emph{data} path are piecewise linear. The
online-observation homotopy of \citet{Garrigues2008} exploits the second: a scalar
$t$ interpolates a rank-one change in the data,
\begin{equation}
G(t) = G + t\,\sigma\, u u^\top,
\qquad
c(t) = c + t\,\sigma\, u\, w,
\qquad t \in [0,1],
\label{eq:online_homotopy}
\end{equation}
with $\sigma = +1$ for an entering observation $(u,w)$ and $\sigma = -1$ for a
leaving one. Along $t$ the active coefficients move linearly until a breakpoint ---
an active coordinate reaching zero, or an inactive correlation reaching $\pm\mu_j$
--- at which the active set pivots and a new segment begins. The breakpoints are
roots of scalar rational functions of $t$ and are finite in number, so the walk
from $t=0$ to $t=1$ \emph{terminates exactly}: no convergence tolerance, no
approximate fallback. One window slide is therefore exactly two calls
to~\eqref{eq:online_homotopy}.

Three implementation choices make this usable. \emph{Per-coordinate penalties:}
$\mu$ is a vector, zero on the intercept and on the dense HAR block, and those
coordinates are additionally \emph{locked} into the active set and never tested for
exit --- an elastic net on the exogenous columns wrapped around an unpenalized
least-squares fit of HAR, expressed in the optimizer rather than approximated by a
small $\alpha$. \emph{Warm starting:} the active set and signs carry from bar to
bar, so a typical slide costs a few pivots, not a path walk from $\mu_{\max}$.
\emph{Periodic cold reseed:} $G$ and $c$ are themselves maintained by rank-one
add/remove, so an exact rebuild from the current window block periodically
re-anchors the state and bounds float drift.

The positioning is the conceptual content here. For ridge, a rank-one data change
produces a rank-one change in the inverse and~\eqref{eq:sherman_morrison} delivers
the new solution in one algebraic step. For the lasso no such fixed formula exists,
because the identity of the active set is itself data-dependent. The online
homotopy \emph{is} the $L_1$ analogue of Sherman--Morrison: same role, same
rank-one trigger, but a variable-cost piecewise walk in place of a fixed-cost
identity --- which is also why $L_1$ needs the periodic exact refresh ridge does
not. The verdict (Section~\ref{sec:linvsnonlin}): on the widest bucket the causally
tuned recursive elastic net scores $0.12827$ against ridge's $0.12788$, a gap of
$0.0004$ beside the $0.006$ by which either beats the incumbent. The penalty family
is not a lever; the algorithm makes that comparison exact rather than
cadence-confounded.

\subsubsection{Numerical honesty: the identifiability mask}
\label{sec:alg_mask}

Exact every-bar updating has an uncomfortable corollary. At $\alpha = 0$ nothing
regularizes the design, and the frozen feature set contains blocks that are
singular by construction on particular windows: a source going live part-way
through the sample has availability indicators that are identically zero before its
go-live bar and identically one after, so a window entirely on one side sees them
as constant and a window straddling it sees several of them collapse onto
\emph{the same} step function.

The fix is a per-window identifiability mask, recomputed from the raw training
window at every refit, in three passes ordered cheapest first. Pass~1 drops
columns with no dispersion ($\max = \min$); pass~2 drops exact byte-duplicates of
an earlier survivor; both are threshold-free and scale-free. Pass~3 handles what
pairwise comparison cannot see --- a block that is rank-deficient without any two
columns being equal. On the survivors, centred and rescaled to unit column norm
with singular values $\sigma_1 \ge \cdots \ge \sigma_m$, it takes the rank cut
\begin{equation}
\mathrm{rank} = \#\Bigl\{\, i \;:\; \sigma_i > \max(n,m)\, \varepsilon_{\mathrm{mach}}\, \sigma_1 \,\Bigr\},
\label{eq:rank_cut}
\end{equation}
exactly the tolerance \texttt{numpy.linalg.matrix\_rank} uses --- machine precision
relative to the design's own largest singular value, on a centred unit-norm
design, so neither the cut nor the pivot ordering is an artefact of feature units.
The retained set is the column-pivoted QR pivot order truncated at that rank:
greedy most-independent-first, so the dropped columns are the redundant ones.
Masked columns are \emph{dropped}, never zeroed: a zeroed column leaves $X^\top X$
singular and \texttt{Ridge(alpha=0)} then returns coefficients of order $10^{13}$
on it, harmless only for as long as the column is exactly zero. The failure this
fixes was measured, not anticipated: on the \texttt{implied\_vol} bucket at one
local window the constant/duplicate passes kept $45$ columns whose rank was $41$,
and the unpenalized fit predicted $2.2\times10^{12}$ where neighbouring bars
predicted $1.3$. Two of $100$ walk-forward chunks blew up this way, and because the
fitted residual series feeds downstream stages each poisoned every remaining bar of
its chunk: roughly $2\%$ of bars carried the pooled QLIKE from $0.13$ to $1.12$.

\subsection{Causal hyperparameter re-selection}
\label{sec:alg_tune}

The second lever is hyperparameter choice, and the protocol is periodic and
strictly causal. Every $250$ solves the current training window of length $W$ is
split \emph{forward} into a fit block, an embargo, and a validation tail:
\begin{equation}
\underbrace{[\,t-W,\; t-V-E\,)}_{\text{fit}}
\quad
\underbrace{[\,t-V-E,\; t-V\,)}_{\text{embargo } E = 25}
\quad
\underbrace{[\,t-V,\; t\,)}_{\text{validation } V = 125}.
\label{eq:forward_split}
\end{equation}
Every candidate on a log-scale penalty grid is fitted on the fit block and scored
by one-step-ahead mean squared error on the validation tail; the argmin holds until
the next re-selection. Only rows strictly inside the training window are ever read,
so the region $[\,t,\infty)$ is untouched. Forward validation rather than
cross-validation is deliberate: $K$-fold assumes a within-window exchangeability
that volatility regimes violate, and it is not the deployment structure. The
embargo covers the $25$-bar HAR moving average, which would otherwise bleed
fit-block information into validation through the feature construction itself;
longer HAR lags cannot be embargoed inside a $500$-bar window, which we state
rather than paper over. Between re-selections the solution slides exactly:
ridge through~\eqref{eq:sherman_morrison}, the $L_1$ arms through two calls
to~\eqref{eq:online_homotopy}. Each re-selection is also a cold reseed on the full
window, doubling as the exact re-anchor that bounds float drift, and recomputes the
identifiability mask of Section~\ref{sec:alg_mask} from the raw ring buffer.

The verdict is negative and is reported in full in
Section~\ref{sec:linvsnonlin}: causally tuned arms \emph{trail} their fixed-default
counterparts. That is the point of building the protocol. Without it, ``we tried
tuning and it did not help'' is an anecdote about one analyst's grid; with it, the
statement is that on this panel, under a protocol in which every forecast's
hyperparameters were selected from its own past, adaptive penalty selection is not
a source of skill.

\subsection{Amortized composition}
\label{sec:alg_amortized}

The residual-stack models of Section~\ref{sec:linvsnonlin} are two-stage: a linear
base $\hat{g}_t$ on the training window, and a gradient-boosted tree ensemble
\citep{Chen2016,Ke2017} on the base's residuals $y - \hat{g}$. Tuning the tree
stage means sweeping tree configurations, and done naively each trial re-runs the
whole composition --- including a per-bar \texttt{sklearn} \texttt{Ridge} base,
$231{,}000$ fits per trial, roughly $825$ seconds per evaluation even on the
smallest cell. That puts a battery out of reach. The base, however, is
\emph{config-independent within a cell}: it depends on the bucket, the training
window and the base penalty, none of which the tree sweep varies. So it is computed
once with the rank-one solver and cached --- both the every-bar out-of-sample base
predictions and the base coefficients at each tree-refit point $t_r$, the latter
being what forms the in-sample residual $y_{\mathrm{train}} -
X_{\mathrm{train}}\beta_{t_r}$ that the tree at that refit trains on. Each trial
then fits only trees, and the composed forecast is
\begin{equation}
\hat{y}_t \;=\; \hat{g}_t \;+\; \mathcal{T}_{t_r}(x_t),
\qquad t_r = \text{the refit block containing } t,
\label{eq:amortized}
\end{equation}
which reproduces the naive composed backtest bit for bit. That claim is not left
to argument: the implementation carries a \texttt{selfcheck} mode whose only job
is to run both paths and assert equality. Amortization here is an exactness
property, not an approximation, and it is what made the tree side of the
Section~\ref{sec:linvsnonlin} battery feasible.

\subsection{Spectral \texorpdfstring{$k$}{k}-NN: nonparametric regression via manifold learning}
\label{sec:alg_spectral}

The three algorithms above accelerate, discipline, or compose a parametric fit. The
fourth changes the geometry. Its research question is one volatility forecasters
ask and rarely test: \emph{do realized-variance dynamics recur?} If the market
passes through a limited repertoire of states, then for any bar there exist
\emph{analog episodes} in the past --- stretches whose recent trajectory and
contemporaneous conditions resemble the present --- and what happened next in those
episodes is informative about what happens next now. A global parametric model
cannot express that; a nearest-neighbour regression can, provided the notion of
resemblance is right, and manifold learning is the natural candidate for supplying
it.

\subsubsection{View construction and the global normalization}

The unit of comparison is a \emph{view}: a vector for bar $j$ concatenating a
\textbf{path} block --- the trailing $W$-bar trajectory of the base model's
residual $e$ --- with a \textbf{state} block --- the non-target-derived feature
columns $x_j$ of the bar the view predicts:
\begin{equation}
v_j \;=\; \Bigl[\;
\underbrace{\tfrac{e_{j-W:j}}{\mathrm{IQR}_w \cdot \sqrt{W+d}}}_{\text{path, } W \text{ dims}}
\;,\;
\underbrace{\tfrac{x_j}{\sqrt{W+d}}}_{\text{state, } d \text{ dims}}
\;\Bigr] \;\in\; \mathbb{R}^{W+d}.
\label{eq:view}
\end{equation}
Four properties are decisions, not defaults. \emph{The path is uncentred:} a common
level cancels in a difference of views, so centring would erase exactly what
distinguishes a calm episode from a violent one; leaving it in means genuine level
gaps register, which is the regime-analog property the estimator is supposed to
have. \emph{The path scale is one causal scalar per refit:} $\mathrm{IQR}_w$ is the
interquartile range of the training window's own target values, read only off rows
already handed to the model, which puts the unscaled path on the same dimensionless
footing as the state block, itself rolling-robust-scaled to unit IQR upstream; a
degenerate window ($\mathrm{IQR}_w \le 0$) \emph{zeroes} the path block rather than
dividing by a floor, so there is no epsilon in~\eqref{eq:view}. \emph{The state
block excludes target-derived columns:} the HAR moving averages and their calendar
interactions are dropped, because the path already carries the target's history at
full resolution; re-admitting them would re-encode that information and silently
re-weight it in the distance --- severely for narrow buckets, negligibly for wide
ones --- confounding the bucket axis with ``how much does the state duplicate the
path''. The stated cost is that the view's reach into the past is exactly $W$ bars.
\emph{Selection is by column name}, off the harness's own feature names, and the
constructor raises rather than silently re-including the target block.

The single scalar $\sqrt{W+d}$ divides \emph{both} blocks, so equal weight falls on
each \emph{coordinate} and the state block's share of squared distance is
$d/(W+d)$: $0.9\%$ for the empty-exog \texttt{baseline} bucket ($d = 9$ calendar
dimensions), $35\%$ for \texttt{all\_features} ($d = 525$). The share scales with
how much information the bucket carries, with no tuning constant. The obvious
alternative --- normalize each block by its own dimension --- was implemented
first, measured, and rejected. It equalizes the two \emph{blocks}, handing a
nine-dimensional one-hot calendar state the same total distance weight as a
$960$-dimensional volatility path; and a one-hot encoding is categorical, so
day-of-week identity outweighed the volatility path and \emph{cut the $k$-NN
graph}. Measured at the first refit of the first task (\texttt{baseline},
$23{,}040$ views, graph $k = 10$), per-block normalization produced three connected
components of sizes $14{,}016 / 4{,}512 / 4{,}512$, the two smaller of which were
$100\%$ single-day-of-week (purity $1.000$). This is fatal rather than cosmetic: on
a disconnected graph the Laplacian's bottom eigenvectors \emph{are} the component
indicators, so the ``embedding'' encodes the day of the week and nothing about
temporal geometry --- and \texttt{sklearn} warned exactly that at every refit. One
global $\sqrt{W+d}$ collapses this to a single component at day-of-week purity
$0.197 \approx 1/5$, i.e.\ chance. Raising the graph $k$ also reconnects the graph
but leaves the distorted metric intact, so it was rejected as treating the symptom;
a smooth cyclical clock encoding is the principled way to retain the session-edge
channel and is left to its own arm.

\subsubsection{Embedding, extension, regression, and the OLS anchor}

Given the $n \times (W+d)$ matrix of training views, a binary $k$-nearest-neighbour
graph is built and the embedding is the bottom $d_\phi$ nontrivial eigenvectors of
the normalized graph Laplacian --- Laplacian eigenmaps \citep{BelkinNiyogi2003}.
Spectral embeddings have no native \texttt{transform}: they are defined only on
the points they were fitted on. Placing a new bar in the learned coordinate system
requires an out-of-sample extension \citep{Bengio2004}, here a Gaussian-weighted
$k$-NN Nystr\"om interpolation. For a test view $v$ with graph neighbours $\{i\}$
at distances $\delta_i$,
\begin{equation}
\phi(v) = \frac{\sum_i w_i \, \phi_i}{\sum_i w_i},
\qquad
w_i = \exp\!\Bigl(-\frac{\delta_i^2}{2\sigma^2}\Bigr),
\qquad
\sigma = \operatorname{median}_i \delta_i,
\label{eq:nystrom}
\end{equation}
the bandwidth set per query by the median neighbour distance, so it adapts to local
density rather than being a hyperparameter. The regression stage is a $k$-NN
regressor in embedding space under the same self-tuning kernel; the forecast
composes as
\begin{equation}
\hat{y}_t \;=\; \underbrace{\hat{g}_t(x_t)}_{\text{OLS base}}
\;+\; \omega \cdot \underbrace{\hat{c}_t\bigl(\phi(v_t)\bigr)}_{\text{$k$-NN correction on the base's residuals}},
\qquad \omega \in (0,1],
\label{eq:spectral_pred}
\end{equation}
with $\omega = 1$ throughout. Setting $d_\phi = 0$ bypasses the spectral projection
and runs the $k$-NN directly on the composed view of~\eqref{eq:view}; we call this
the \emph{passthrough} arm, and it is the ablation that isolates what manifold
learning contributes.

The anchor $\hat{g}_t$ was not in the original design. The estimator was first run
\emph{unanchored} --- an identity residualizer, so the path block is the target
itself and the $k$-NN forecasts $y$ from scratch --- and came back badly
miscalibrated: on the canary chunk the Mincer--Zarnowitz slope was $0.474$, with
QLIKE $0.391$ against the incumbent's $0.179$ on that chunk and
$R^2_{\mathrm{oos}} = -2.23$. A slope near one-half with a large negative $R^2$ is
the signature of over-shrinkage toward the sample mean: local averaging estimates a
conditional mean that has been flattened, because the neighbourhoods are not tight
enough to resolve the level. The remedy is to stop asking the analog search to
reproduce HAR's level from scratch. The residualizer is set to exact OLS --- the
production ridge path at $\alpha = 0$, routed through the identifiability mask
of~\eqref{eq:rank_cut} --- which is precisely the estimator defining the paper's
incumbent, so the model reads as \emph{incumbent plus local correction}. Everything
else is held fixed, leaving the identity arm a clean control.

\subsubsection{Computational notes}

The eigenproblem is the bottleneck, and became one only after the graph was fixed.
On the real first-refit matrix ($23{,}040$ views $\times\ 969$ dimensions, graph
$k = 10$), an LOBPCG eigensolver with an algebraic-multigrid preconditioner took
$19.8$ s locally against ARPACK's $>22$ minutes on a faster cluster node --- better
than $60\times$, with finite and sanely scaled output. Why ARPACK had been fast
\emph{before} is itself the point: while the graph was disconnected, the eigenvalue
$0$ had multiplicity three and the component indicators \emph{were} the
eigenvectors, so the problem was trivial; connecting the graph made it genuinely
hard, and the solver had to change with it. A different eigensolver returns the
same invariant subspace but not bit-identical eigenvectors, so downstream neighbour
sets can shift slightly. Cadences differ by cost: the residualizer refits every bar
--- closed-form, cheap, and the anchor --- while the spectral basis and the $k$-NN
regressor rebuild every $240$ bars, with intervening bars placed
by~\eqref{eq:nystrom}. Thread pools are sized from the scheduler's stated
entitlement rather than the node's core count, the direct fix for the BLAS
allocation failures over-subscription caused.

\subsubsection{Results}

% Source: writeup/metrics_table_causal_tune_plus_spectral.csv (family = spectral-kNN),
% all 32 arms, verbatim. Incumbent QLIKE 0.13415 and n = 218,934 from that file's
% incumbent_qlike / n columns. Run health and rejections: writeup/experiment_ledger.csv.
\begin{table}[H]
\centering
\footnotesize
\setlength{\tabcolsep}{3pt}
\caption{Spectral $k$-NN arms on the $n = 218{,}934$ panel, all against the same
incumbent (per-bar OLS on HAR $+$ calendar, QLIKE $0.13415$). $W$ = view window in
bars, $d_\phi$ = embedding dimension ($0 =$ passthrough), $k$ = regression
neighbours. Negative DM favours the spectral arm. The nine production-default arms
are identically configured and differ only in bucket; the best and worst are shown
and the remaining seven lie between them, spanning DM $+45.99$ to $+52.38$.}
\label{tab:spectral_arms}
\begin{tabular}{lrrrllrrr}
\toprule
Bucket & $W$ & $d_\phi$ & $k$ & Anchor & Mask & QLIKE & DM & $p$ \\
\midrule
\multicolumn{9}{l}{\textit{Panel A --- beats the incumbent}} \\
market\_vw   & 12 & 0 & 100 & OLS & const$+$dup & 0.13096 & $-17.57$ & $4.4\!\times\!10^{-69}$ \\
market\_vw   & 12 & 0 & 100 & OLS & ---         & 0.13145 & $-14.60$ & $2.9\!\times\!10^{-48}$ \\
baseline     & 24 & 0 & 100 & OLS & ---         & 0.13218 & $-12.91$ & $4.2\!\times\!10^{-38}$ \\
baseline     & 12 & 0 & 100 & OLS & ---         & 0.13219 & $-14.37$ & $7.7\!\times\!10^{-47}$ \\
baseline     & 12 & 6 & 100 & OLS & ---         & 0.13364 & $-4.61$  & $4.0\!\times\!10^{-6}$ \\
\midrule
\multicolumn{9}{l}{\textit{Panel B --- statistically indistinguishable}} \\
baseline     & 24 & 0 & 25  & OLS & ---         & 0.13397 & $-0.88$ & $0.378$ \\
baseline     & 24 & 6 & 100 & OLS & ---         & 0.13399 & $-1.61$ & $0.108$ \\
baseline     & 12 & 0 & 25  & OLS & ---         & 0.13405 & $-0.52$ & $0.603$ \\
\midrule
\multicolumn{9}{l}{\textit{Panel C --- loses}} \\
baseline     & 12 & 6 & 25  & OLS      & ---    & 0.13676 & $13.16$ & $1.5\!\times\!10^{-39}$ \\
baseline     & 24 & 6 & 25  & OLS      & ---    & 0.13701 & $15.66$ & $2.7\!\times\!10^{-55}$ \\
baseline     & 12 & 0 & 25  & identity & ---    & 0.14173 & $19.12$ & $1.6\!\times\!10^{-81}$ \\
baseline     & 12 & 0 & 100 & identity & ---    & 0.14236 & $18.28$ & $1.2\!\times\!10^{-74}$ \\
baseline     & 24 & 0 & 25  & identity & ---    & 0.14475 & $21.17$ & $1.6\!\times\!10^{-99}$ \\
baseline     & 24 & 0 & 100 & identity & ---    & 0.14539 & $20.79$ & $5.4\!\times\!10^{-96}$ \\
implied\_vol & 12 & 0 & 25  & identity & ---    & 0.14938 & $28.35$ & $9.1\!\times\!10^{-177}$ \\
sentiment    & 12 & 0 & 25  & identity & ---    & 0.15214 & $31.40$ & $2.2\!\times\!10^{-216}$ \\
market\_vw   & 12 & 0 & 25  & identity & ---    & 0.15773 & $37.93$ & $9.4\!\times\!10^{-315}$ \\
baseline     & 12 & 6 & 25  & identity & ---    & 0.15993 & $45.71$ & $<\!10^{-300}$ \\
baseline     & 12 & 6 & 100 & identity & ---    & 0.16213 & $41.22$ & $<\!10^{-300}$ \\
vol\_demand  & 12 & 0 & 25  & identity & ---    & 0.16490 & $28.38$ & $3.8\!\times\!10^{-177}$ \\
baseline     & 24 & 6 & 25  & identity & ---    & 0.16872 & $47.08$ & $<\!10^{-300}$ \\
baseline     & 24 & 6 & 100 & identity & ---    & 0.16874 & $40.91$ & $<\!10^{-300}$ \\
all\_features& 12 & 0 & 25  & identity & ---    & 0.17409 & $37.79$ & $1.9\!\times\!10^{-312}$ \\
\addlinespace
\multicolumn{9}{l}{\textit{Original production default --- all nine buckets, range over buckets}} \\
best (moments)    & 960 & 8 & 25 & identity & --- & 0.24627 & $51.16$ & $<\!10^{-300}$ \\
worst (all\_feat.) & 960 & 8 & 25 & identity & --- & 0.26353 & $43.35$ & $<\!10^{-300}$ \\
\bottomrule
\end{tabular}
\end{table}

Table~\ref{tab:spectral_arms} reports every scored arm. The best is OLS-anchored,
masked, short-window ($W = 12$), \emph{passthrough} ($d_\phi = 0$), heavily
averaged ($k = 100$), on the \texttt{market\_vw} bucket: QLIKE $0.13096$ against
$0.13415$, DM $-17.57$, $p = 4.4\times10^{-69}$; its unmasked sibling reaches
$0.13145$ at DM $-14.60$, and the corresponding empty-exog arms land at
$0.13218$--$0.13219$. A middle group --- the $k = 25$ OLS-anchored arms --- is
indistinguishable from the incumbent ($p$ between $0.108$ and $0.603$). The losing
group is large and ordered: every unanchored arm loses by DM $+13$ to $+38$, and
the original production default loses catastrophically on all nine buckets, QLIKE
$0.246$ to $0.264$ at DM $+43$ to $+52$ --- the configuration one would have
shipped from the module's defaults is roughly twice as bad as the incumbent. The
$16$ \texttt{baseline}-bucket arms form a complete $2^4$ factorial in $W$,
$d_\phi$, $k$ and anchor, which separates the levers cleanly
(Table~\ref{tab:spectral_effects}).

% Source: writeup/metrics_table_causal_tune_plus_spectral.csv, the 16 arm_tag="fac_*"
% rows (bucket = baseline), a complete 2^4 factorial. Main effects are differences of
% 8-cell means; conditional effects are means of the 4 matched pairs.
\begin{table}[H]
\centering
\small
\caption{Main and conditional effects on QLIKE from the complete $2^4$ factorial on
the \texttt{baseline} bucket ($16$ arms, $n = 218{,}934$ each). Positive is
\emph{worse}.}
\label{tab:spectral_effects}
\begin{tabular}{llrl}
\toprule
Lever & Contrast & Main effect & Conditional on anchor \\
\midrule
Anchor      & identity $-$ OLS          & $+0.02000$ & --- \\
Embedding   & $d_\phi{=}6 - d_\phi{=}0$ & $+0.01179$ & OLS $+0.00225$; identity $+0.02132$ \\
View window & $W{=}24 - W{=}12$         & $+0.00274$ & OLS $+0.00012$; identity $+0.00536$ \\
Neighbours  & $k{=}100 - k{=}25$        & $-0.00079$ & OLS $-0.00245$; identity $+0.00087$ \\
\bottomrule
\end{tabular}
\end{table}

\paragraph{Finding 1: the OLS anchor is load-bearing.}
Replacing exact OLS by the identity costs $+0.02000$ QLIKE averaged over the
factorial --- the largest effect in the grid, nearly twice the embedding effect and
seven times the window effect. At the otherwise-best matched cell ($W = 12$,
$d_\phi = 0$, $k = 100$) the two arms are $0.14236$ against $0.13219$, a gap of
$0.01017$; across the eight matched pairs the gap ranges from $0.00768$ to
$0.03475$, widest where the other settings are worst. Read with the calibration
diagnostics above, this says the analog search cannot recover the conditional level
of realized variance on its own. Given the level for free it has something to
contribute; asked to produce the level, it does not.

\paragraph{Finding 2: the spectral embedding hurts everywhere.}
Passthrough beats the $d_\phi = 6$ embedding in \emph{every} matched cell:
$+0.00145$, $+0.00181$, $+0.00271$ and $+0.00304$ on the four OLS-anchored pairs,
$+0.0182$ to $+0.0240$ on the four unanchored ones. This is the section's most
uncomfortable result --- manifold learning is the part of the design that was
supposed to be novel --- so the mechanism deserves an explanation rather than a
shrug. Two independent measurements supply one. First, the out-of-sample extension
is imprecise relative to the structure it must resolve: holding out $462$ days,
fitting the embedding on the rest, extending the held-out days
via~\eqref{eq:nystrom} and comparing to a full-sample embedding under a Procrustes
alignment (spectral coordinates being defined only up to rotation and sign), the
median per-point placement error is $74.4\%$ of the embedding's own scale ---
$0.0037$ against a coordinate scale of $0.0049$. A test point cannot be placed on
the manifold at the precision the manifold's geometry demands. Second, the manifold
holds little geometric information to begin with. In the prototype study the
leading coordinate $\phi_1$ tracks \texttt{har\_ma\_1} while calendar years
intermingle --- the structure is a one-dimensional volatility continuum, not a set
of discrete regimes, and named crises smear along the high-amplitude tail rather
than forming separated clusters. Over $4{,}127$ walk-forward test days the spectral
$k$-NN scored $R^2 = 0.408$ against $0.419$ for a direct $k$-NN in the raw feature
space, the two prediction series correlating $0.96$, and both were beaten by simple
persistence ($R^2 = 0.496$). If $\phi_1$ is essentially the volatility level, a
bar's $\phi$-neighbours are its raw-space neighbours and the spectral step is a
lossy re-parameterization of a neighbourhood the $k$-NN already had. Consistently,
a diagnostic on the first production refit found neighbours in the raw composed
space had targets $15.0\%$ closer than random against only $10.3\%$ after
projection onto eight Laplacian eigenvectors --- the projection discards about a
third of the analog signal, even though the projected space's distance ratios look
dramatically better, which is a low-dimensionality artefact rather than predictive
value.

\paragraph{Finding 3: large-$k$ local averaging is what survives.}
Among OLS-anchored arms, moving from $k = 25$ to $k = 100$ improves QLIKE in all
four matched pairs, by $0.00179$ to $0.00312$ (mean $0.00245$), and the
distinction is decisive: every $k = 100$ anchored arm sits in Panel~A of
Table~\ref{tab:spectral_arms}, every $k = 25$ anchored passthrough arm in Panel~B.
Small-$k$ analog retrieval --- the version of the idea that ``finds the few most
similar past episodes'' --- injects variance the weak per-neighbour signal cannot
pay for. Note the interaction: on unanchored arms the ordering reverses slightly
(mean $+0.00087$), which is why the \emph{main} effect of $k$ is small and why
reading main effects alone would have hidden the result. Averaging over a hundred
neighbours is not analog retrieval in any strong sense; it is local smoothing.

\paragraph{What actually survives measurement.}
The honest summary is that the manifold-learning component --- the part that
motivated the whole construction --- is the part that does not pay. The embedding
hurts in every matched cell, its out-of-sample extension is imprecise by
three-quarters of the embedding's own scale, and the manifold it learns is a
one-dimensional volatility continuum the raw distance already encodes. What remains
is an OLS-anchored short-window analog search with heavy local averaging: a
nonparametric correction to a parametric forecast, in the raw composed view space,
with $k$ large enough that the correction is a smoother rather than a retriever.
That object is genuinely good --- it beats the OLS incumbent decisively, $0.13096$
against $0.13415$ at DM $-17.57$, on $218{,}934$ bars against a shared
per-bar-refit comparator --- and, just as clearly, not competitive with the model
classes of Section~\ref{sec:linvsnonlin}, where the best causally tuned tree
reaches $0.12604$ and the best causally tuned linear arm $0.12788$. Its correct
placement is as a \emph{different} kind of improvement over the incumbent, not a
better one, and as evidence that the recurrence hypothesis in the form tested here
is worth less than reading the exogenous panel with a conventional estimator.

\paragraph{Robustness caveats.}
Three limitations are material. \emph{First, the winning bucket is the only
surviving one.} Of the exogenous buckets run at the winning settings,
\texttt{market\_vw} is the only one whose arm completed cleanly; its siblings were
rejected as singular --- \texttt{implied\_vol}, \texttt{sentiment} and
\texttt{vol\_demand} each blew two of one hundred walk-forward chunks and
\texttt{all\_features} blew seven, with pooled QLIKE of $1.06$--$3.11$ and
Mincer--Zarnowitz slopes between $10^{-24}$ and $10^{-29}$, the unmistakable
signature of the mechanism of Section~\ref{sec:alg_mask}. That \texttt{market\_vw}
wins may therefore reflect which bucket happened to survive rather than which
carries the most analog signal, so the exogenous axis is one clean observation, not
a ranking. \emph{Second, the mask in the results table is the partial one.} The
arms marked \texttt{const}$+$\texttt{dup} used passes~1 and~2 only; the
rank-revealing pass~3 of~\eqref{eq:rank_cut} was written \emph{in response} to the
failures those arms exposed, and no fully re-run pooled arm under the complete
three-pass mask is reported here. Correspondingly no \texttt{all\_features}
OLS-anchored spectral number is valid: every attempt blew chunks, and the only one
producing a non-degenerate pooled figure --- a masked $d_\phi = 32$ variant ---
pooled just $35$ of $100$ chunks with seven blown, so it is recorded in the run
ledger and excluded from the table. \emph{Third, the composition weight is
unexplored.} Every arm runs at $\omega = 1$ in~\eqref{eq:spectral_pred}. Given that
the estimator's dominant failure mode is over-confidence of the nonparametric
stage, shrinking $\omega$ is the most plausible remaining improvement, and it has
not been measured; the one attempt at $\omega = 0.25$ produced no usable output.

\subsection{Synthesis}
\label{sec:alg_synthesis}

% Sources by row: src/models/rolling_least_squares.py, src/backtest/rolling_ridge.py
% (timing from writeup/session_notes_2026-06-23_feature_investigation.md);
% src/models/reclasso_har.py; specs/causal_tune_linear.py; resid_amortized.py;
% src/models/spectral_knn.py + src/features/extractors/spectral_embedding.py with
% writeup/metrics_table_causal_tune_plus_spectral.csv.
\begin{table}[H]
\centering
\footnotesize
\setlength{\tabcolsep}{4pt}
\caption{The four algorithm designs: idea, device, verdict. Every verdict is on the
$n = 218{,}934$ panel except the timings, which are per-bar wall-clock
measurements.}
\label{tab:alg_synthesis}
\begin{tabular}{p{2.2cm} p{3.3cm} p{3.4cm} p{4.6cm}}
\toprule
Algorithm & Mathematical idea & Computational device & Verdict \\
\midrule
Rank-one rolling LS
& Sliding-window sufficient statistics change by rank one; $\alpha = 0$ is exact OLS
& Two rank-one updates per slide, $O(p^2)$ not $O(Wp^2)$; Sherman--Morrison on the ridged inverse for the tuned variant
& Exact to $1.8\times10^{-11}$ vs.\ \texttt{sklearn}, $108\times$ faster ($4.589 \to 0.042$ ms/bar). This \emph{is} the incumbent, and what makes per-bar refitting affordable throughout. \\
\addlinespace
Online enet homotopy
& Elastic net as lasso on a ridged Gram; exact piecewise-linear $L_1$ path in the rank-one data parameter $t \in [0,1]$
& Exactly two homotopy calls per bar; finite breakpoints give exact termination with no fallback; warm active set, periodic cold reseed
& The $L_1$ analogue of Sherman--Morrison. Matches ridge to $0.0004$ QLIKE: the penalty family is a non-lever, and the algorithm makes that comparison exact rather than cadence-confounded. \\
\addlinespace
Causal rolling tuning
& Every forecast's hyperparameters chosen only from its own past: fit / embargo / forward-validation split inside the window
& Grid scored on the forward split every $250$ solves; exact rank-one sliding between tunings; cold reseed doubles as drift re-anchor
& Causally tuned arms \emph{trail} fixed defaults. Tuning is a non-lever --- and the protocol converts that from anecdote to causal claim. \\
\addlinespace
Amortized residual trees
& The linear base is config-independent within a cell, so it need not be recomputed per trial
& Precompute every-bar base predictions and per-refit base coefficients once; trials fit trees only on $y - \hat{g}$; \texttt{selfcheck} asserts bit-equality
& Bit-exact by construction and by assertion; removes $231{,}000$ ridge refits and $\approx 825$ s per evaluation. Enabled the tree side of the Section~\ref{sec:linvsnonlin} battery. \\
\addlinespace
Spectral $k$-NN
& Analog retrieval: nonparametric regression in a metric learned by Laplacian eigenmaps over path$+$state views
& One global $\sqrt{W+d}$ normalization (a measured fix for a graph cut along day-of-week); AMG-preconditioned LOBPCG, $19.8$ s vs.\ $>22$ min; Nystr\"om extension between $240$-bar refits
& Beats the incumbent decisively ($0.13096$ vs.\ $0.13415$, DM $-17.57$) --- but only once reduced to an OLS-anchored, passthrough, large-$k$ local average. The embedding is rejected by measurement in every matched cell, and the arm still trails every causally tuned arm of Section~\ref{sec:linvsnonlin}. \\
\bottomrule
\end{tabular}
\end{table}

Read as a whole, the section says where the remaining leverage does \emph{not} lie.
Exactness costs nothing once the right update is derived, but it is a statement
about computation, not skill: the rank-one solver and the online homotopy make the
incumbent and its penalized variants cheap and honest without changing what they
know. Causal tuning is worth having for the same reason and with the same limit.
Amortization buys a battery, not a forecast. And the one genuinely different
geometry tested here --- learn a manifold, retrieve analogs on it --- beats the
incumbent only in the degenerate limit where the manifold has been removed and
retrieval has become smoothing. Under a frozen feature set, algorithm design bought
exactness, feasibility, and the defensibility of the negative results; it did not
buy the level of the forecast. Consistently with Sections~\ref{sec:marginal}
and~\ref{sec:linvsnonlin}, the lever is information.
